% Options for packages loaded elsewhere
\PassOptionsToPackage{unicode}{hyperref}
\PassOptionsToPackage{hyphens}{url}
%
\documentclass[
]{article}
\usepackage{amsmath,amssymb}
\usepackage{iftex}
\ifPDFTeX
  \usepackage[T1]{fontenc}
  \usepackage[utf8]{inputenc}
  \usepackage{textcomp} % provide euro and other symbols
\else % if luatex or xetex
  \usepackage{unicode-math} % this also loads fontspec
  \defaultfontfeatures{Scale=MatchLowercase}
  \defaultfontfeatures[\rmfamily]{Ligatures=TeX,Scale=1}
\fi
\usepackage{lmodern}
\ifPDFTeX\else
  % xetex/luatex font selection
\fi
% Use upquote if available, for straight quotes in verbatim environments
\IfFileExists{upquote.sty}{\usepackage{upquote}}{}
\IfFileExists{microtype.sty}{% use microtype if available
  \usepackage[]{microtype}
  \UseMicrotypeSet[protrusion]{basicmath} % disable protrusion for tt fonts
}{}
\makeatletter
\@ifundefined{KOMAClassName}{% if non-KOMA class
  \IfFileExists{parskip.sty}{%
    \usepackage{parskip}
  }{% else
    \setlength{\parindent}{0pt}
    \setlength{\parskip}{6pt plus 2pt minus 1pt}}
}{% if KOMA class
  \KOMAoptions{parskip=half}}
\makeatother
\usepackage{xcolor}
\usepackage{listings}
\newcommand{\passthrough}[1]{#1}
\lstset{defaultdialect=[5.3]Lua}
\lstset{defaultdialect=[x86masm]Assembler}
\usepackage{longtable,booktabs,array}
\usepackage{calc} % for calculating minipage widths
% Correct order of tables after \paragraph or \subparagraph
\usepackage{etoolbox}
\makeatletter
\patchcmd\longtable{\par}{\if@noskipsec\mbox{}\fi\par}{}{}
\makeatother
% Allow footnotes in longtable head/foot
\IfFileExists{footnotehyper.sty}{\usepackage{footnotehyper}}{\usepackage{footnote}}
\makesavenoteenv{longtable}
\setlength{\emergencystretch}{3em} % prevent overfull lines
\providecommand{\tightlist}{%
  \setlength{\itemsep}{0pt}\setlength{\parskip}{0pt}}
\setcounter{secnumdepth}{5}
\ifLuaTeX
  \usepackage{selnolig}  % disable illegal ligatures
\fi
\usepackage{bookmark}
\IfFileExists{xurl.sty}{\usepackage{xurl}}{} % add URL line breaks if available
\urlstyle{same}
\hypersetup{
  pdftitle={Orin Cluster - Current Architecture and Operations},
  pdfauthor={Samuel Padilla},
  hidelinks,
  pdfcreator={LaTeX via pandoc}}

\title{Orin Cluster - Current Architecture and Operations}
\usepackage{etoolbox}
\makeatletter
\providecommand{\subtitle}[1]{% add subtitle to \maketitle
  \apptocmd{\@title}{\par {\large #1 \par}}{}{}
}
\makeatother
\subtitle{Current-state technical overview and operations reference}
\author{Samuel Padilla}
\date{June 4, 2026}

\begin{document}
\maketitle

{
\setcounter{tocdepth}{3}
\tableofcontents
}
\section{Purpose}\label{purpose}

Orin Cluster is a Kubernetes-based local edge-AI assistant platform. The
active architecture is centered on \textbf{Cortex} as the local ingress,
routing, command, and orchestration process. Runtime work is represented
with canonical \textbf{WorkPackets} and routed to discovered backend
services.

The current implementation is no longer best described as a
gateway/orchestrator/Ollama split. The active shape is:

\begin{itemize}
\tightlist
\item
  \textbf{Cortex}: ingress, Harness, command interface, deployment
  operations, discovery, routing, local fast response, JobManager, and
  Memory/session integration.
\item
  \textbf{LLM}: Primer-backed WorkPacket executor for model inference.
\item
  \textbf{Memory}: WorkPacket backend for users, sessions, chat history,
  claims, notes, summaries, prompt context, and Federation records.
\item
  \textbf{Tool}: planned backend role for non-LLM tools. The role is
  represented in protocols and deployment paths, but no complete tool
  backend is documented here yet.
\item
  \textbf{Federation}: external communication boundary for public
  network discovery, token-gated joining, signed WorkPacket envelopes,
  capability advertisement, and federated work/result exchange.
\item
  \textbf{Common}: shared protocol, runtime, Primer, adapter,
  JobManager, model-loading, and Hugging Face helpers.
\end{itemize}

The current direction is:

\begin{lstlisting}
WorkPacket is the common execution unit.
RouterService is the common routing/submission path.
JobManager is the common async job and pipeline mechanism.
Federation is the external membrane and must not expose internal cluster structure.
\end{lstlisting}

\begin{center}\rule{0.5\linewidth}{0.5pt}\end{center}

\section{Repository structure}\label{repository-structure}

Current project shape:

\begin{lstlisting}
orin-cluster/
  canonical-unified-image/
  documentation/
  src/
    common/
    cortex/
    federation/
    llm/
    memory/
    tool/
  install-orin.sh
  terminal-chat.py
\end{lstlisting}

Important source areas:

\begin{lstlisting}
src/common/
  Shared protocols, RuntimeMessage, WorkPacket, WorkResult, Primer,
  adapters, JobManager, and runtime policies.

src/cortex/
  Cortex app, Harness, command router, deployment service,
  Kubernetes discovery, routing, planner, and backend client.

src/llm/
  Primer-backed LLM backend app.

src/memory/
  Memory backend app, SQLite/file-backed persistence, session/message logic,
  prompt context, and Federation storage operations.

src/federation/
  Federation models, settings, Memory-backed storage, network registry,
  join tokens, members, identity/signing, envelopes, policy,
  capability summaries, Cortex bridge, and outbound client.

src/tool/
  Planned backend role for non-LLM tools.
\end{lstlisting}

\begin{center}\rule{0.5\linewidth}{0.5pt}\end{center}

\section{Image model}\label{image-model}

The intended image split is:

\begin{lstlisting}
standard runtime, amd64/arm64:
  FROM orin-gw:5000/primer-runtime:0.2

jetson runtime:
  FROM orin-gw:5000/primer-runtime:jetson-orin.0.2
\end{lstlisting}

The application folders can have standard and Jetson Dockerfiles:

\begin{lstlisting}
Dockerfile
Dockerfile.jetson
\end{lstlisting}

The long-term model should distinguish only between:

\begin{lstlisting}
standard runtime: amd64 and arm64
Jetson runtime: Jetson-specific base/runtime
\end{lstlisting}

Current cleanup note: some deployment code has carried three-way
platform image selection for \passthrough{\lstinline!jetson!},
\passthrough{\lstinline!arm64!}, and \passthrough{\lstinline!amd64!}.
The preferred direction is for \passthrough{\lstinline!amd64!} and
\passthrough{\lstinline!arm64!} to share the standard runtime image
naming while Jetson keeps its separate runtime image.

\begin{center}\rule{0.5\linewidth}{0.5pt}\end{center}

\section{\texorpdfstring{Bootstrap model:
\texttt{install-orin.sh}}{Bootstrap model: install-orin.sh}}\label{bootstrap-model-install-orinsh}

\passthrough{\lstinline!install-orin.sh!} is the initial host/bootstrap
installer. It is not the long-term general pod deployment mechanism.

Its responsibilities are:

\begin{enumerate}
\def\labelenumi{\arabic{enumi}.}
\tightlist
\item
  detect host platform and GPU state
\item
  optionally configure the local test registry for k3s/containerd
\item
  install k3s server on the bootstrap/control host
\item
  create namespace and RBAC
\item
  create the Cortex SSH-key secret
\item
  create the k3s join-token secret
\item
  write the deployment configuration file
\item
  create the deployment configuration ConfigMap
\item
  label the bootstrap node
\item
  deploy the initial Memory app
\item
  wait for Memory rollout
\item
  deploy the initial Cortex app
\item
  wait for Cortex rollout
\end{enumerate}

\subsection{4.1 Default bootstrap
values}\label{41-default-bootstrap-values}

Important defaults:

\begin{lstlisting}
ORIN_NAMESPACE=orin
HOST_ALIAS=$(hostname)
HOST_NAME=$(hostname)
HOST_IP=$(hostname -I | awk '{print $1}')
K3S_NODE_NAME=$HOST_ALIAS
SSH_USER=ubuntu
SSH_KEY=~/.ssh/id_ed25519
SSH_KEY_MOUNT_PATH=/data/ssh/id_ed25519
SSH_KEY_SECRET_NAME=cortex-ssh-key
SSH_PORT=22
PLATFORM=auto
NVIDIA_GPU=auto
CORTEX_SIZE=light
CORTEX_NODE_PORT=30080
CORTEX_SERVICE_ACCOUNT=cortex
REGISTRY=orin-gw:5000
IMAGE_VERSION=1.0.0
CONFIG_PATH=/data/configuration
K3S_DATA_DIR=/data/k3s
K3S_STORAGE_DIR=/data/k3s-storage
CONFIGURE_LOCAL_REGISTRY=true
LOCAL_REGISTRY_ENDPOINT=http://${REGISTRY}
\end{lstlisting}

\subsection{4.2 Configuration file}\label{42-configuration-file}

The installer writes JSON configuration to:

\begin{lstlisting}
/data/configuration
\end{lstlisting}

The Cortex pod receives that file through the ConfigMap:

\begin{lstlisting}
deployment-configuration
\end{lstlisting}

mounted at:

\begin{lstlisting}
/data/configuration
\end{lstlisting}

The configuration seeds the initial host/control node and deployment
defaults. It is not the final source of truth for every future node.
Additional nodes are added later through Cortex\textquotesingle s
command interface and \passthrough{\lstinline!DeploymentService!}.

\subsection{4.3 Secrets}\label{43-secrets}

The installer creates:

\begin{lstlisting}
cortex-ssh-key
  mounted in Cortex at /data/ssh/id_ed25519

k3s-join-token
  contains key: token
\end{lstlisting}

Cortex uses the SSH key for node operations and reads the k3s join token
through the Kubernetes API when installing k3s agents on additional
nodes.

\subsection{4.4 Bootstrap deployments}\label{44-bootstrap-deployments}

The installer deploys the minimum working runtime:

\begin{lstlisting}
memory
cortex
\end{lstlisting}

Cortex is exposed as:

\begin{lstlisting}
Service: cortex-service
Type: NodePort
NodePort: 30080
\end{lstlisting}

Memory is exposed internally as:

\begin{lstlisting}
Service: memory-service
Type: ClusterIP
\end{lstlisting}

Federation deployment support exists as a target direction, but this
document treats the initial installer as a Cortex/Memory bootstrap path
unless the installer has explicitly been extended for Federation in the
current branch.

\subsection{4.5 Bootstrap RBAC}\label{45-bootstrap-rbac}

The Cortex service account needs permissions to manage cluster resources
and discover backends.

The installer creates a \passthrough{\lstinline!ClusterRole!} and
\passthrough{\lstinline!ClusterRoleBinding!} for:

\begin{lstlisting}
nodes: get, list, watch, patch, update
namespaces: get, list, create
pods: get, list, watch, create, patch, update, delete
services: get, list, watch, create, patch, update, delete
persistentvolumeclaims: get, list, watch, create, patch, update, delete
configmaps: get, list, watch, create, patch, update, delete
secrets: get, list, watch
deployments: get, list, watch, create, patch, update, delete
endpointslices: get, list, watch
\end{lstlisting}

EndpointSlice access is required because backend discovery counts ready
endpoints from EndpointSlices.

\begin{center}\rule{0.5\linewidth}{0.5pt}\end{center}

\section{Runtime applications}\label{runtime-applications}

\subsection{5.1 Cortex}\label{51-cortex}

Cortex is the main local runtime. It owns:

\begin{itemize}
\tightlist
\item
  public ingress for terminal/chat requests
\item
  Harness response flow
\item
  slash-command handling
\item
  local Primer-backed fast response path
\item
  user/session/message persistence through Memory WorkPackets
\item
  backend discovery and registry refresh
\item
  backend routing and WorkPacket submission
\item
  deployment operations through
  \passthrough{\lstinline!DeploymentService!}
\item
  JobManager-based background jobs and pipelines
\item
  direct WorkPacket handling for internal/Federation use
\end{itemize}

At startup, Cortex wires components such as:

\begin{lstlisting}
Primer
DeploymentService
DiscoveryService
BackendRoutingPolicy
BackendClient
Planner
RouterService
CommandRouter
JobManager
Harness
\end{lstlisting}

Cortex starts discovery refresh during lifespan startup and stops the
Primer backend on shutdown.

\subsection{5.2 Harness}\label{52-harness}

Harness is the ingress layer for terminal and chat requests. It
validates sessions, persists user messages through Memory, performs
lightweight interpretation, optionally executes deterministic terminal
actions, optionally starts background pipelines, and returns direct or
fast responses.

Harness supports two ingress modes:

\begin{lstlisting}
terminal
chat
\end{lstlisting}

Terminal mode is command/help/status oriented and should not start
background jobs.

Chat mode normally answers directly. It starts background work only when
the lightweight interpreter chooses the
\passthrough{\lstinline!start\_pipeline!} action.

\subsection{5.3 LLM}\label{53-llm}

The LLM app is a Primer-backed WorkPacket executor. It exposes
backend/model lifecycle endpoints and a packet execution API.

It supports:

\begin{lstlisting}
POST /engine or equivalent backend-engine loading path
POST /load
POST /unload
POST /work
GET  /work/{work_id}
GET  /model_status
GET  /live
GET  /ready
GET  /health
GET  /attach
\end{lstlisting}

The active discovery convention expects model/runtime status at:

\begin{lstlisting}
/model_status
\end{lstlisting}

This replaces older documentation that treated
\passthrough{\lstinline!/models!} as the active discovery status
endpoint.

Cortex discovers LLM services through Kubernetes service metadata and
routes \passthrough{\lstinline!work\_type=llm!} packets to them.

\subsection{5.4 Memory}\label{54-memory}

The Memory app is a synchronous WorkPacket backend. Cortex and
Federation route Memory WorkPackets to it for:

\begin{itemize}
\tightlist
\item
  user upsert
\item
  session resolution
\item
  chat message creation
\item
  recent conversation retrieval
\item
  claims
\item
  notes
\item
  summaries
\item
  prompt context
\item
  Federation persistent records
\end{itemize}

The Memory backend normally returns completed
\passthrough{\lstinline!WorkResult!} values directly from
\passthrough{\lstinline!POST /work!}.

\subsection{5.5 Tool}\label{55-tool}

The Tool role is present in protocols, routing hints, deployment
commands, and Federation public capability mapping. A complete current
tool backend is not documented here yet.

\subsection{5.6 Federation}\label{56-federation}

Federation is the external communication boundary for Orin clusters.

It handles:

\begin{lstlisting}
public network discovery
token-gated network joining
member identity
signed member requests
signed federated work envelopes
policy enforcement
sanitized capability advertisement
federated work submission
federated result polling
\end{lstlisting}

Federation must not expose internal Orin cluster structure. External
peers must not see Cortex URLs, LLM URLs, Memory URLs, Kubernetes
services, pod names, node IPs, SSH details, deployment configuration,
model paths, memory contents, slash commands, backend load/unload
controls, or internal backend refs.

External communication goes through Federation endpoints only.
Federation validates, verifies, authorizes, sanitizes, and then hands
accepted work into local Cortex through
\passthrough{\lstinline!CortexBridge!}.

\begin{center}\rule{0.5\linewidth}{0.5pt}\end{center}

\section{Runtime endpoints}\label{runtime-endpoints}

\subsection{6.1 Cortex endpoints}\label{61-cortex-endpoints}

\begin{longtable}[]{@{}lrl@{}}
\toprule\noalign{}
Endpoint & Method & Purpose \\
\midrule\noalign{}
\endhead
\bottomrule\noalign{}
\endlastfoot
\passthrough{\lstinline!/work!} & POST & Main ingress. In chat/terminal
flows this is session-based; internal/Federation work can use
WorkPacket-shaped requests where supported by the current Cortex app. \\
\passthrough{\lstinline!/work/\{work\_id\}!} & GET & Poll
internal/deferred work result where supported. \\
\passthrough{\lstinline!/network!} & GET & List discovered backend
descriptors. \\
\passthrough{\lstinline!/network/\{role\}!} & GET & List discovered
candidates for a role. \\
\passthrough{\lstinline!/load!} & POST & Start local Cortex model-load
job when using Cortex-local Primer. \\
\passthrough{\lstinline!/unload!} & POST & Stop local Cortex Primer
backend/model. \\
\passthrough{\lstinline!/jobs/\{job\_id\}!} or command
\passthrough{\lstinline!/job\_status!} & GET/command & Inspect jobs
depending on app surface. \\
\passthrough{\lstinline!/attach!} & GET & Stream local Cortex logs when
configured. \\
\passthrough{\lstinline!/live!} & GET & Process liveness. \\
\passthrough{\lstinline!/ready!} & GET & Readiness and Primer status. \\
\passthrough{\lstinline!/health!} & GET & Basic health endpoint. \\
\end{longtable}

\subsection{6.2 LLM endpoints}\label{62-llm-endpoints}

\begin{longtable}[]{@{}lrl@{}}
\toprule\noalign{}
Endpoint & Method & Purpose \\
\midrule\noalign{}
\endhead
\bottomrule\noalign{}
\endlastfoot
\passthrough{\lstinline!/engine!} or backend-load equivalent & POST &
Load/select \passthrough{\lstinline!gguf!} or
\passthrough{\lstinline!vllm!} backend engine. \\
\passthrough{\lstinline!/load!} & POST & Start model-load job. \\
\passthrough{\lstinline!/jobs/\{job\_id\}!} & GET & Read backend job
status. \\
\passthrough{\lstinline!/unload!} & POST & Stop current
backend/model. \\
\passthrough{\lstinline!/work!} & POST & Accept a
\passthrough{\lstinline!WorkPacket!}. \\
\passthrough{\lstinline!/work/\{work\_id\}!} & GET & Poll submitted LLM
work. \\
\passthrough{\lstinline!/model\_status!} & GET & Return loaded
model/runtime metadata, including effective context. \\
\passthrough{\lstinline!/attach!} & GET & Stream backend logs. \\
\passthrough{\lstinline!/live!} & GET & Process liveness. \\
\passthrough{\lstinline!/ready!} & GET & Readiness and Primer status. \\
\passthrough{\lstinline!/health!} & GET & Basic health endpoint. \\
\end{longtable}

\subsection{6.3 Memory endpoints}\label{63-memory-endpoints}

\begin{longtable}[]{@{}lrl@{}}
\toprule\noalign{}
Endpoint & Method & Purpose \\
\midrule\noalign{}
\endhead
\bottomrule\noalign{}
\endlastfoot
\passthrough{\lstinline!/work!} & POST & Execute Memory
\passthrough{\lstinline!WorkPacket!} operations synchronously. \\
\passthrough{\lstinline!/work/\{work\_id\}!} & GET & Placeholder or
lookup path; Memory work usually completes on POST. \\
\passthrough{\lstinline!/live!} & GET & Process liveness. \\
\passthrough{\lstinline!/ready!} & GET & Readiness. \\
\passthrough{\lstinline!/health!} & GET & Basic health endpoint. \\
\end{longtable}

\subsection{6.4 Federation endpoints}\label{64-federation-endpoints}

\begin{longtable}[]{@{}lrl@{}}
\toprule\noalign{}
Endpoint & Method & Purpose \\
\midrule\noalign{}
\endhead
\bottomrule\noalign{}
\endlastfoot
\passthrough{\lstinline!/live!} & GET & Process liveness. \\
\passthrough{\lstinline!/ready!} & GET & Readiness, cluster ID, protocol
version. \\
\passthrough{\lstinline!/health!} & GET & Basic health and cluster
ID. \\
\passthrough{\lstinline!/federation/networks!} & GET & List public
Federation networks. \\
\passthrough{\lstinline!/federation/networks!} & POST & Create local
Federation network. Local/admin endpoint. \\
\passthrough{\lstinline!/federation/networks/\{slug\}!} & GET & Get
public-safe network view. \\
\passthrough{\lstinline!/federation/networks/\{slug\}/tokens!} & POST &
Create join token. Local/admin endpoint. \\
\passthrough{\lstinline!/federation/networks/\{slug\}/join!} & POST &
Join network using token, cluster ID, and public key. \\
\passthrough{\lstinline!/federation/networks/\{slug\}/capabilities!} &
GET & Read public network capabilities. \\
\passthrough{\lstinline!/federation/networks/\{slug\}/capabilities/refresh!}
& POST & Refresh local advertised capabilities from Cortex discovery. \\
\passthrough{\lstinline!/federation/networks/\{network\_id\}/heartbeat!}
& POST & Signed member heartbeat. \\
\passthrough{\lstinline!/federation/networks/\{network\_id\}/capabilities!}
& POST & Signed member capability publish. \\
\passthrough{\lstinline!/federation/networks/\{network\_id\}/work!} &
POST & Submit signed federated WorkPacket envelope. \\
\passthrough{\lstinline!/federation/networks/\{network\_id\}/work/\{work\_id\}!}
& GET & Poll federated work result. \\
\end{longtable}

\begin{center}\rule{0.5\linewidth}{0.5pt}\end{center}

\section{Canonical runtime protocol}\label{canonical-runtime-protocol}

The current architecture uses internal canonical message and work types
instead of treating OpenAI chat messages as the universal internal
contract.

\subsection{7.1 Runtime messages}\label{71-runtime-messages}

\begin{lstlisting}[language=Python]
PartType = Literal["text", "image", "audio", "video", "json", "binary"]
PartEncoding = Literal["plain", "base64"]
Visibility = Literal["user", "internal"]

MessageRole = Literal[
    "system",
    "developer",
    "user",
    "assistant",
    "tool",
    "context",
]

class ContentPart(BaseModel):
    type: PartType
    data: str
    encoding: PartEncoding = "plain"
    mime_type: str | None = None
    metadata: dict[str, Any] = Field(default_factory=dict)

class RuntimeMessage(BaseModel):
    role: MessageRole
    parts: list[ContentPart] = Field(default_factory=list)
    name: str | None = None
    metadata: dict[str, Any] = Field(default_factory=dict)
\end{lstlisting}

Current implementation note:
\passthrough{\lstinline!ContentPart.mime\_type!} may be auto-filled from
\passthrough{\lstinline!type!} and \passthrough{\lstinline!encoding!}
when missing, for example \passthrough{\lstinline!text/plain!}.

When JSON-shaped payloads are embedded inside a text prompt for the
model, prefer \passthrough{\lstinline!type="text"!},
\passthrough{\lstinline!encoding="plain"!}, and
\passthrough{\lstinline!mime\_type="text/plain"!} unless the part is
intended to be treated as an actual
\passthrough{\lstinline!application/json!} content part.

\subsection{7.2 Ingress and egress}\label{72-ingress-and-egress}

\begin{lstlisting}[language=Python]
class InferenceSession(BaseModel):
    user_id: str
    session_id: str | None = None
    channel: str = "api"
    stream: bool = False
    content: list[RuntimeMessage]
    metadata: dict[str, Any]

class EgressResponse(BaseModel):
    content: list[RuntimeMessage]
    work_id: str | None = None
    session_id: str | None = None
    metadata: dict[str, object] = {}
\end{lstlisting}

\subsection{7.3 Work packets}\label{73-work-packets}

\begin{lstlisting}[language=Python]
class WorkType(str, Enum):
    LLM = "llm"
    TOOL = "tool"
    CORTEX = "cortex"
    MEMORY = "memory"

class WorkDisposition(str, Enum):
    DIRECT = "direct"
    DEFERRED = "deferred"

class RoutingHints(BaseModel):
    role: WorkType | None = None
    required_capabilities: list[str] = Field(default_factory=list)
    required_modalities: list[str] = Field(default_factory=list)
    runtime_preference: list[dict[str, Any]] = Field(default_factory=list)
    latency_preference: str = "normal"

class CanonicalTask(BaseModel):
    work_type: WorkType
    operation: str
    messages: list[RuntimeMessage] = Field(default_factory=list)
    memory_request: BaseRuntimeMemoryRequest | None = None
    inputs: dict[str, Any] = Field(default_factory=dict)
    constraints: dict[str, Any] = Field(default_factory=dict)
    routing_hints: RoutingHints = Field(default_factory=RoutingHints)

class WorkPacket(BaseModel):
    work_id: str
    disposition: WorkDisposition = WorkDisposition.DIRECT
    task: CanonicalTask
    metadata: dict[str, Any] = Field(default_factory=dict)

class WorkResult(BaseModel):
    status: str
    work_id: str
    content: list[RuntimeMessage] = Field(default_factory=list)
    backend_name: str | None = None
    backend_model: str | None = None
    error: str | None = None
    metadata: dict[str, Any] = Field(default_factory=dict)
\end{lstlisting}

Allowed \passthrough{\lstinline!WorkResult.status!} values:

\begin{lstlisting}
accepted
running
completed
failed
\end{lstlisting}

Important current field name:

\begin{lstlisting}
memory_request
\end{lstlisting}

Older text or code that says \passthrough{\lstinline!memoryrequest!} is
stale.

\subsection{7.4 WorkResult federation
sanitization}\label{74-workresult-federation-sanitization}

\passthrough{\lstinline!WorkResult.sanitize\_for\_federation()!} returns
a reduced WorkResult intended for external Federation egress.

It preserves:

\begin{lstlisting}
status
work_id
content
backend_name
backend_model
error
\end{lstlisting}

It drops most metadata and only keeps sanitized
\passthrough{\lstinline!work\_details!} when present:

\begin{lstlisting}
work_type
operation
\end{lstlisting}

This prevents internal routing metadata such as backend refs, selected
backend details, trusted internal flags, or local execution context from
leaking through Federation.

\begin{center}\rule{0.5\linewidth}{0.5pt}\end{center}

\section{Harness ingress and response
flow}\label{harness-ingress-and-response-flow}

Harness validates the incoming session, persists the user message
through Memory, performs lightweight interpretation, optionally executes
a deterministic terminal action or starts a background pipeline, and
then returns either a direct response or a fast model-generated
response.

\subsection{8.1 Session validation}\label{81-session-validation}

For each ingress request, Harness:

\begin{enumerate}
\def\labelenumi{\arabic{enumi}.}
\tightlist
\item
  validates the incoming \passthrough{\lstinline!InferenceSession!} as
  an \passthrough{\lstinline!InferenceObject!}
\item
  upserts the user through Memory
\item
  resolves or creates a session through Memory if no
  \passthrough{\lstinline!session\_id!} was provided
\item
  inserts the incoming user message into Memory
\item
  continues to terminal or chat response handling
\end{enumerate}

\subsection{8.2 Lightweight
interpretation}\label{82-lightweight-interpretation}

Harness builds a small interpretation context containing:

\begin{lstlisting}
Ingress context:
- mode
- user_id
- session_id
- channel

Session background jobs:
- compact job summary, when jobs exist

Available mode commands/actions:
- command_router.get_commands(mode)
\end{lstlisting}

This context is intentionally small. It gives the lightweight
interpreter enough information to choose an action without exposing full
backend state or large runtime details.

\subsection{8.3 Terminal mode}\label{83-terminal-mode}

Terminal mode is for CLI help, deterministic inspection, and safe
terminal/system information actions. It should not start background
work.

Terminal interpretation model:

\begin{lstlisting}
{
  "mode": "terminal",
  "intent": "terminal_help | system_status | job_status | normal_chat | unknown",
  "action": "answer_only | terminal_info_action | report_job_status",
  "terminal_action": "string | null"
}
\end{lstlisting}

If the action is \passthrough{\lstinline!terminal\_info\_action!},
Harness executes the selected command-router action through:

\begin{lstlisting}[language=Python]
command_router.execute_harness_action(
    terminal_action,
    mode="terminal",
)
\end{lstlisting}

Only safe information actions are executed automatically.

If the action is \passthrough{\lstinline!report\_job\_status!}, Harness
reads session jobs from JobManager and formats session job status.

\subsection{8.4 Chat mode}\label{84-chat-mode}

Chat mode is the normal assistant path.

Chat interpretation model:

\begin{lstlisting}
{
  "mode": "chat",
  "action": "answer_only | start_pipeline | report_job_status",
  "pipeline": "chat | null"
}
\end{lstlisting}

Normal questions should use \passthrough{\lstinline!answer\_only!}.

Background work starts when the user explicitly asks for deeper,
detailed, researched, verified, current, multi-step, or comprehensive
work. When the action is \passthrough{\lstinline!start\_pipeline!},
Harness starts a JobManager pipeline:

\begin{lstlisting}
kind: harness.chat
pipeline: chat
\end{lstlisting}

The immediate response does not wait for the full result:

\begin{lstlisting}
A deeper response is being worked on. Background job `<job_id>` has been started. Ask for the status or result later.
\end{lstlisting}

If the action is \passthrough{\lstinline!report\_job\_status!}, Harness
returns formatted job status for the current session.

\subsection{8.5 Fast response path}\label{85-fast-response-path}

If Harness does not return a direct terminal/action/job-status response,
it builds a fast response prompt.

The fast response prompt can include:

\begin{itemize}
\tightlist
\item
  current user message
\item
  last assistant message from Memory
\item
  compact response context
\item
  internal response instruction derived from lightweight interpretation
\item
  background job id, if one was started
\end{itemize}

The fast response is generated through the local Primer using the
\passthrough{\lstinline!chat!} generation policy.

\subsection{8.6 Chat background
pipeline}\label{86-chat-background-pipeline}

The current \passthrough{\lstinline!chat!} pipeline stages are:

\begin{lstlisting}
get_prompt
interpret_turn
shape_turn
build_packets
read_and_send_packets
build_final_response_message
\end{lstlisting}

Stage responsibilities:

\begin{lstlisting}
get_prompt:
  Fetch prompt context from Memory for the current session.

interpret_turn:
  Build turn-interpretation messages and route an llm.inspect WorkPacket
  through RouterService. Primer is used for token estimation; actual work
  is routed to an LLM backend.

shape_turn:
  Read the interpretation result, build task-shaping messages, and route
  another llm.inspect WorkPacket through RouterService.

build_packets:
  Convert shaped tasks into canonical WorkPackets.

read_and_send_packets:
  Submit built WorkPackets through router.send_many(), wait for send batch,
  call router.retrieve_many(), and wait for retrieve batch.

build_final_response_message:
  Build the final response prompt from the original user message, prompt
  context, and collated packet results. Then route an llm.analyze WorkPacket
  and poll until completion or failure.
\end{lstlisting}

The shaper output is normalized before packet creation. Accepted shapes
include:

\begin{lstlisting}
task_items
items
tasks
\end{lstlisting}

The normalization step guards against the model incorrectly routing
LLM-style work to Cortex. Operations such as
\passthrough{\lstinline!chat!}, \passthrough{\lstinline!summarize!},
\passthrough{\lstinline!classify!}, \passthrough{\lstinline!extract!},
and \passthrough{\lstinline!analyze!} are normalized toward
\passthrough{\lstinline!llm!} routing when appropriate.

\subsection{8.7 WorkPacket ingress
pipeline}\label{87-workpacket-ingress-pipeline}

Harness also defines a smaller \passthrough{\lstinline!work\_pipeline!}
for direct WorkPacket ingress.

Current behavior:

\begin{enumerate}
\def\labelenumi{\arabic{enumi}.}
\tightlist
\item
  If \passthrough{\lstinline!packet.task.work\_type == "llm"!}, start
  pipeline \passthrough{\lstinline!work\_pipeline!}.
\item
  Return
  \passthrough{\lstinline!WorkResult(status="accepted", work\_id=<job\_id>)!}.
\item
  Store the original packet work id in metadata as
  \passthrough{\lstinline!original\_work\_id!}.
\item
  Poll the JobManager job through
  \passthrough{\lstinline!retrieve\_work\_packet()!}.
\item
  Return a sanitized WorkResult for federation-safe egress.
\end{enumerate}

The current direct WorkPacket path only accepts LLM work. Other work
types return failed with \passthrough{\lstinline!unknown work type!}.

\subsection{8.8 Current implementation
note}\label{88-current-implementation-note}

The Harness implementation contains branches for a
\passthrough{\lstinline!get\_job\_result!} action, but the current
lightweight interpretation model does not include
\passthrough{\lstinline!get\_job\_result!} as a valid action. Treat this
as an implementation stub until the type model is updated.

\begin{center}\rule{0.5\linewidth}{0.5pt}\end{center}

\section{Command interface}\label{command-interface}

The command router exposes a folder-style terminal interface. Commands
and folders are represented as structured objects, and the same command
inventory is also exposed to Harness for lightweight interpretation.

\subsection{9.1 Command model}\label{91-command-model}

Main concepts:

\begin{lstlisting}
ShellState
  path: list[str]
  selected_backend_name: str | None

CommandContext
  router: CommandRouter
  backend: BackendDescriptor | None
  backend_client -> router.backend_client

Command
  command: str
  desc: str
  kind: "text" | "stream"
  handler
  modes: set["terminal" | "chat"]
  harness_action: str | None
  safe_info_action: bool
  suggest_only: bool

CommandFolder
  name: str
  desc: str
  commands: list[Command]
  folders: dict[str, CommandFolder]
  dynamic: bool
\end{lstlisting}

All text command handlers use:

\begin{lstlisting}[language=Python]
handler(ctx: CommandContext, args: str)
\end{lstlisting}

Stream commands use the same context but return an async iterator.

\passthrough{\lstinline!CommandContext!} is important because
backend-scoped commands no longer need special hardcoded router paths.
If the shell is currently inside
\passthrough{\lstinline!Root / Cluster / <backend>!}, the context
contains the selected backend. Otherwise
\passthrough{\lstinline!ctx.backend!} is \passthrough{\lstinline!None!}.

\subsection{9.2 Global commands}\label{92-global-commands}

Global commands:

\begin{lstlisting}
/help
/commands
/cd
/render
/attach
\end{lstlisting}

\passthrough{\lstinline!/attach!} is a stream command. In a backend
folder it attaches to the selected backend\textquotesingle s
\passthrough{\lstinline!/attach!} endpoint. Outside a backend folder it
attaches to the local Cortex log stream, if configured.

\subsection{9.3 Current folder tree}\label{93-current-folder-tree}

\begin{lstlisting}
Root
  /help
  /commands
  /cd
  /attach
  /render

  Deployment
    /list_pods
    /list_nodes
    /deploy_pod
    /delete_pod
    /install_k3s
    /add_node
    /remove_node
    /label_node
    /update_discovery_meta

  Cluster
    /show
    /cluster_snapshot
    <dynamic backend folders>

  Configuration
    /model_aliases
    /job_status
    /load_model
    /unload_model
    /engine

    ModelDownloader
      /huggingface
\end{lstlisting}

The old backend \passthrough{\lstinline!/network!} command has been
replaced by \passthrough{\lstinline!/show!} in the Cluster context.

\subsection{9.4 Dynamic Cluster folder}\label{94-dynamic-cluster-folder}

\passthrough{\lstinline!Cluster!} is a dynamic folder. When the user
enters:

\begin{lstlisting}
/cd Cluster
\end{lstlisting}

the router lists discovered backends from the backend registry and
exposes each backend as a child folder.

Example:

\begin{lstlisting}
Root / Cluster
  /show
  /cluster_snapshot

  llm-medium-1-service
  memory-service
\end{lstlisting}

When the user enters a backend folder:

\begin{lstlisting}
/cd llm-medium-1-service
\end{lstlisting}

the router sets:

\begin{lstlisting}[language=Python]
shell_state.path = ["Root", "Cluster", backend.name]
shell_state.selected_backend_name = backend.name
\end{lstlisting}

The current folder is then resolved from the selected backend role. If
the backend role is known, the role-specific folder is used. Otherwise
the generic backend folder is used.

\subsection{9.5 Backend command
folders}\label{95-backend-command-folders}

Generic backend commands:

\begin{lstlisting}
/show
/health
/update_discovery_meta
\end{lstlisting}

LLM backend commands:

\begin{lstlisting}
/show
/health
/update_discovery_meta
/model_status
/load_model
/job_status
/unload_model
\end{lstlisting}

Memory currently uses the generic backend command folder.

Future role folders can be added for Tool, Cortex, and Federation.

\subsection{9.6 Human visibility versus Harness action
inventory}\label{96-human-visibility-versus-harness-action-inventory}

The human terminal remains folder-oriented.
\passthrough{\lstinline!/commands!} shows the commands visible in the
current folder plus global commands.

Harness is different. The AI should not be limited by the human
user\textquotesingle s current CLI folder.
\passthrough{\lstinline!CommandRouter.get\_commands(mode)!} returns the
broader action inventory available to the lightweight interpreter.

Each returned command entry contains:

\begin{lstlisting}
action
command
description
folder
safe_info_action
suggest_only
requires_backend
backend_role
\end{lstlisting}

This lets Harness decide whether a user request can be fulfilled by a
safe deterministic terminal action.

\subsection{9.7 Safe and suggest-only
actions}\label{97-safe-and-suggest-only-actions}

Commands are tagged for Harness use.

\begin{lstlisting}
safe_info_action=True
\end{lstlisting}

means Harness may execute the command automatically for information
retrieval.

\begin{lstlisting}
suggest_only=True
\end{lstlisting}

means Harness should not execute the command automatically. It can
suggest the command or explain it, but execution requires the user to
explicitly issue that command.

Examples of safe information actions:

\begin{lstlisting}
list_commands
render_current_folder
list_pods
list_nodes
show_cluster
cluster_snapshot
model_status
job_status
show_backend
backend_health
\end{lstlisting}

Examples of suggest-only or mutating actions:

\begin{lstlisting}
change_folder
attach_logs
deploy_pod
delete_pod
install_k3s
add_node
remove_node
label_node
update_discovery_meta
load_model
unload_model
engine_status
huggingface
\end{lstlisting}

\passthrough{\lstinline!execute\_harness\_action()!} enforces this at
runtime. It rejects unknown actions, non-safe actions, and suggest-only
actions.

\subsection{9.8 Model and Hugging Face
commands}\label{98-model-and-hugging-face-commands}

Current model alias examples:

\begin{lstlisting}
qwen35-0.8b
qwen35-2B
qwen35-4b
qwen35-9b
\end{lstlisting}

The \passthrough{\lstinline!/huggingface!} command supports:

\begin{lstlisting}
/huggingface search <query>
/huggingface info <repo_id>
/huggingface gguf <repo_id>
/huggingface alias add <alias> <repo_id> [--file filename] [--tokenizer tokenizer_id]
\end{lstlisting}

\begin{center}\rule{0.5\linewidth}{0.5pt}\end{center}

\section{JobManager}\label{jobmanager}

\passthrough{\lstinline!JobManager!} is the shared runtime mechanism for
asynchronous jobs, background pipelines, command-triggered long-running
work, router batch work, and Harness background work.

It is generic and is not model-load specific.

\subsection{10.1 Core job types}\label{101-core-job-types}

\begin{lstlisting}[language=Python]
JobStatus = Literal["accepted", "running", "completed", "failed"]
StageStatus = Literal["accepted", "running", "completed", "failed"]

@dataclass(slots=True)
class JobSpec:
    kind: str
    payload: dict[str, Any] = field(default_factory=dict)
    job_id: str | None = None

@dataclass(slots=True)
class Job:
    job_id: str
    spec: JobSpec
    status: JobStatus = "accepted"
    error: str | None = None

    pipeline_name: str | None = None
    current_stage: str | None = None
    stage_index: int = 0
    stage_count: int = 0
    stage_status: StageStatus | None = None
    stage_results: dict[str, dict[str, Any]] = field(default_factory=dict)
    latest_result: dict[str, Any] = field(default_factory=dict)
    progress_message: str | None = None
\end{lstlisting}

A job can be either a single runner job or a pipeline job with multiple
stages.

Single jobs are used for simple async operations such as router
send/retrieve jobs and local model-load jobs.

Pipeline jobs are used for staged workflows such as Harness background
chat.

\subsection{10.2 Pipeline model}\label{102-pipeline-model}

A pipeline is a named list of stages:

\begin{lstlisting}[language=Python]
@dataclass(slots=True)
class PipelineStage:
    name: str
    runner: PipelineStageRunner
    poller: PipelineStagePoller | None = None
    timeout_seconds: float = 120.0
    poll_interval_seconds: float = 1.0

@dataclass(slots=True)
class PipelineDefinition:
    name: str
    stages: list[PipelineStage]
\end{lstlisting}

If the runner returns \passthrough{\lstinline!completed!} or
\passthrough{\lstinline!failed!}, the stage ends immediately.

If the runner returns \passthrough{\lstinline!accepted!} or
\passthrough{\lstinline!running!}, the stage must have a poller. The
poller is called until the result becomes
\passthrough{\lstinline!completed!} or \passthrough{\lstinline!failed!},
or until the stage timeout expires.

If a stage returns \passthrough{\lstinline!accepted!} or
\passthrough{\lstinline!running!} without a poller, the stage fails.

\subsection{10.3 Separate job and pipeline
semaphores}\label{103-separate-job-and-pipeline-semaphores}

\passthrough{\lstinline!JobManager!} uses two separate semaphores:

\begin{lstlisting}
job_semaphore
pipeline_semaphore
\end{lstlisting}

This avoids a deadlock where a parent pipeline holds the only job slot
while waiting for child router jobs that also need the same slot.

Normal jobs run under \passthrough{\lstinline!job\_semaphore!}.

Pipeline jobs run under \passthrough{\lstinline!pipeline\_semaphore!}.

This means a running pipeline can safely start child jobs through
RouterService and wait for their batch results.

\subsection{10.4 Starting jobs and
pipelines}\label{104-starting-jobs-and-pipelines}

Normal job:

\begin{lstlisting}[language=Python]
job_manager.start(
    spec=JobSpec(kind="...", payload={...}),
    runner=runner,
    on_completed=optional_hook,
)
\end{lstlisting}

Pipeline job:

\begin{lstlisting}[language=Python]
job_manager.define_pipeline("chat", stages=[...])

job_manager.start_pipeline(
    spec=JobSpec(kind="harness.chat", payload={...}),
    pipeline="chat",
)
\end{lstlisting}

The returned job starts as \passthrough{\lstinline!accepted!}. When
execution begins, it becomes \passthrough{\lstinline!running!}.

When a runner returns, the result is converted to a JSON-safe dictionary
with:

\begin{lstlisting}[language=Python]
formatting.as_dict(result)
\end{lstlisting}

and stored in:

\begin{lstlisting}
job.latest_result
\end{lstlisting}

If the runner returns a \passthrough{\lstinline!WorkResult!} with
\passthrough{\lstinline!status="failed"!}, the job is marked failed and
\passthrough{\lstinline!job.error!} is set from the WorkResult error.

\subsection{10.5 Batch work}\label{105-batch-work}

\passthrough{\lstinline!JobManager!} supports lightweight batches
through \passthrough{\lstinline!BatchWork!}:

\begin{lstlisting}[language=Python]
@dataclass(slots=True)
class BatchWork:
    job_ids: list[str] = field(default_factory=list)
    batch_id: str = ""
\end{lstlisting}

A batch stores only job IDs. It does not store
\passthrough{\lstinline!asyncio.Task!} objects.

Creating a batch:

\begin{lstlisting}[language=Python]
batch_id = job_manager.create_batch(jobs)
\end{lstlisting}

Waiting for a batch:

\begin{lstlisting}[language=Python]
await job_manager.batch_progress(batch_id)
\end{lstlisting}

\passthrough{\lstinline!batch\_progress()!} returns a list of normalized
job results once every job in the batch has status
\passthrough{\lstinline!completed!} or \passthrough{\lstinline!failed!}.

Default behavior:

\begin{lstlisting}
blocking: true
poll interval: 1 second
timeout: 300 seconds
\end{lstlisting}

\subsection{10.6 Job inspection}\label{106-job-inspection}

Useful inspection methods:

\begin{lstlisting}
get(job_id)
get_progress(job_id)
list_jobs()
list_active_jobs()
list_jobs_for_session(session_id)
list_active_jobs_for_session(session_id)
summarize_jobs_for_prompt(session_id)
\end{lstlisting}

Current status meanings:

\begin{lstlisting}
accepted  = job has been created and queued
running   = job or pipeline stage is currently running
completed = job finished successfully
failed    = job failed
\end{lstlisting}

Current implementation note:
\passthrough{\lstinline!summarize\_jobs\_for\_prompt()!} includes
payload details for completed jobs. This is useful while debugging but
may be too noisy for model-facing summaries. Long term, completed job
summaries should prefer a compact result summary.

\begin{center}\rule{0.5\linewidth}{0.5pt}\end{center}

\section{Deployment service and
inventory}\label{deployment-service-and-inventory}

\passthrough{\lstinline!DeploymentService!} owns the mutable in-memory
\passthrough{\lstinline!Inventory!}. It reads
\passthrough{\lstinline!/data/configuration!} during startup and seeds
the initial host/control node into inventory.

\subsection{11.1 Inventory model}\label{111-inventory-model}

\begin{lstlisting}
Inventory
  nodes: list[Node]

Node
  alias
  machine_settings
    host
    ip
    platform: amd64 | arm64 | jetson
    nvidia_gpu
  ssh_config
    default_user
    default_key
    port
    auth_method
  kubernetes_settings
    k3s_node_name
    namespace
    data_dir=/data/k3s
    default_local_storage_path=/data/k3s-storage
\end{lstlisting}

\subsection{11.2 DeploymentService
responsibilities}\label{112-deploymentservice-responsibilities}

\passthrough{\lstinline!DeploymentService!} supports:

\begin{itemize}
\tightlist
\item
  adding/removing inventory nodes
\item
  listing inventory nodes with Kubernetes status
\item
  labeling Kubernetes nodes
\item
  checking passwordless sudo over SSH
\item
  installing k3s server/agent over SSH
\item
  deploying pod apps through the Kubernetes Python client
\item
  deleting pod apps
\item
  listing deployed pod apps
\item
  patching service discovery metadata
\item
  resolving pod placement
\end{itemize}

\subsection{11.3 Node placement labels}\label{113-node-placement-labels}

Node labels are used for scheduling and compatibility:

\begin{lstlisting}
orin.role.backend=true
orin.platform=<amd64|arm64|jetson>
orin.gpu=<true|false>
\end{lstlisting}

Optional role preference labels are soft scheduling hints:

\begin{lstlisting}
orin.role.cortex=true
orin.role.llm=true
orin.role.tool=true
orin.role.memory=true
\end{lstlisting}

These are not the same as service discovery labels.

\subsection{11.4 Service discovery
metadata}\label{114-service-discovery-metadata}

Service labels and annotations are used for backend discovery and
routing.

Required service labels:

\begin{lstlisting}
labels:
  orin.ai/backend: "true"
  orin.ai/role: "<role>"
\end{lstlisting}

Supported discovery annotations:

\begin{lstlisting}
annotations:
  orin.ai/kind: "<kind>"
  orin.ai/model: "<model-id>"
  orin.ai/capabilities: "[\"chat\"]"
  orin.ai/modalities: "[\"text\"]"
  orin.ai/priority: "100"
  orin.ai/weight: "1.0"
  orin.ai/visibility: "internal"
\end{lstlisting}

Current convention-based paths:

\begin{lstlisting}
health_path = /health
work_path = /work
model_status_path = /model_status for llm and cortex
model_status_path = None for other roles
\end{lstlisting}

Older annotations such as these are no longer required by the current
discovery provider:

\begin{lstlisting}
orin.ai/health_path
orin.ai/work_path
orin.ai/models_path
\end{lstlisting}

The current discovery provider hardcodes the standard endpoint
convention instead.

\begin{center}\rule{0.5\linewidth}{0.5pt}\end{center}

\section{Pod application model}\label{pod-application-model}

The deployment code uses generic app specs:

\begin{lstlisting}
PodDeploymentSpec
PodServiceSpec
PodPvcSpec
PodDiscoverySpec
PodProbeSpec
PlacementDecision
\end{lstlisting}

For an app named \passthrough{\lstinline!<app\_name>!}, generated
resources are:

\begin{lstlisting}
Deployment: <app_name>
Service:    <app_name>-service
PVC:        <app_name>-pvc
\end{lstlisting}

Cortex is deployed as NodePort. LLM, Memory, Tool, and Federation app
services should generally be ClusterIP unless explicitly exposed.

Current probe convention:

\begin{lstlisting}
readiness: /ready
liveness:  /live
startup:   /ready
\end{lstlisting}

Current storage profiles:

\begin{lstlisting}
light  -> 20Gi
medium -> 50Gi
high   -> 100Gi
\end{lstlisting}

Current cleanup note: some command/type definitions may still use
\passthrough{\lstinline!strong!} instead of
\passthrough{\lstinline!high!}.

\begin{center}\rule{0.5\linewidth}{0.5pt}\end{center}

\section{Discovery, health, and
routing}\label{discovery-health-and-routing}

Cortex discovers routable backends from Kubernetes Services. Discovery
is service-based, not pod-based.

A backend is any Kubernetes Service in the configured Cortex namespace
with:

\begin{lstlisting}
orin.ai/backend=true
\end{lstlisting}

The service must also provide a role label:

\begin{lstlisting}
orin.ai/role=<role>
\end{lstlisting}

Known backend roles:

\begin{lstlisting}
cortex
llm
tool
memory
\end{lstlisting}

\subsection{13.1 Discovery service
composition}\label{131-discovery-service-composition}

\passthrough{\lstinline!DiscoveryService!} wires together:

\begin{lstlisting}
KubernetesDiscoveryProvider
BackendHealthProber
InMemoryBackendRegistry
RegistryRefresher
BackendSelector
\end{lstlisting}

It is initialized with defaults such as:

\begin{lstlisting}
refresh_interval_seconds = 15.0
health_timeout = 2.0
\end{lstlisting}

The registry is refreshed periodically by
\passthrough{\lstinline!RegistryRefresher!}.

\subsection{13.2 Kubernetes service
discovery}\label{132-kubernetes-service-discovery}

\passthrough{\lstinline!KubernetesDiscoveryProvider!} lists Services in
the configured Cortex namespace using:

\begin{lstlisting}
label_selector = "orin.ai/backend=true"
\end{lstlisting}

It also lists EndpointSlices in the same namespace and counts ready
endpoints per service.

A service is skipped if:

\begin{lstlisting}
orin.ai/role is missing
no usable service port exists
\end{lstlisting}

The service port selection prefers a port named
\passthrough{\lstinline!http!}; otherwise discovery uses the first
service port.

\subsection{13.3 Backend URL
convention}\label{133-backend-url-convention}

Backend URLs are built from internal Kubernetes DNS:

\begin{lstlisting}
http://<service-name>.<namespace>.svc.cluster.local:<port>
\end{lstlisting}

Example:

\begin{lstlisting}
http://llm-medium-1-service.orin.svc.cluster.local:8080
\end{lstlisting}

This internal service URL is used only inside the Orin cluster.
Federation must not expose it externally.

\subsection{13.4 BackendDescriptor}\label{134-backenddescriptor}

Discovered services become \passthrough{\lstinline!BackendDescriptor!}
objects:

\begin{lstlisting}[language=Python]
@dataclasses.dataclass(slots=True)
class BackendDescriptor:
    name: str
    namespace: str
    service_name: str
    url: str

    role: BackendRole | str
    kind: str | None = None
    model: str | None = None

    capabilities: list[str] = field(default_factory=list)
    modalities: list[str] = field(default_factory=list)

    priority: int = 100
    weight: float = 1.0
    visibility: str = "internal"

    health_path: str = "/health"
    work_path: str = "/"
    model_status_path: str | None = None

    labels: dict[str, str] = field(default_factory=dict)
    annotations: dict[str, str] = field(default_factory=dict)

    health: BackendHealth = field(default_factory=BackendHealth)
    runtime: RuntimeMetaData = field(default_factory=RuntimeMetaData)
\end{lstlisting}

The descriptor name is currently:

\begin{lstlisting}
<namespace>/<service-name>
\end{lstlisting}

Example:

\begin{lstlisting}
orin/llm-medium-1-service
\end{lstlisting}

\subsection{13.5 Health probing}\label{135-health-probing}

\passthrough{\lstinline!BackendHealthProber!} probes every discovered
backend.

Probe flow:

\begin{enumerate}
\def\labelenumi{\arabic{enumi}.}
\tightlist
\item
  Check whether the service has at least one ready EndpointSlice
  endpoint.
\item
  If no ready endpoints exist, mark backend unavailable.
\item
  If ready endpoints exist, call
  \passthrough{\lstinline!<backend.url>/health!}.
\item
  If \passthrough{\lstinline!/health!} fails, mark backend degraded.
\item
  If the backend has \passthrough{\lstinline!model\_status\_path!}, call
  \passthrough{\lstinline!<backend.url>/model\_status!}.
\item
  If model status fails, mark backend degraded.
\item
  If all required checks pass, mark backend healthy.
\end{enumerate}

For \passthrough{\lstinline!llm!} and \passthrough{\lstinline!cortex!}
backends, discovery attempts model status probing through:

\begin{lstlisting}
/model_status
\end{lstlisting}

If the response contains runtime metadata, the prober updates:

\begin{lstlisting}
backend.runtime.effective_n_ctx
\end{lstlisting}

This value is later used for runtime-aware backend selection.

\subsection{13.6 Runtime metadata}\label{136-runtime-metadata}

Current runtime metadata:

\begin{lstlisting}[language=Python]
@dataclasses.dataclass(slots=True)
class RuntimeMetaData:
    effective_n_ctx: int = 0
\end{lstlisting}

Example runtime preference:

\begin{lstlisting}
runtime_preference = [{"effective_n_ctx": 6000}]
\end{lstlisting}

This means the backend must have
\passthrough{\lstinline!effective\_n\_ctx >= 6000!}.

\subsection{13.7 Capability summary}\label{137-capability-summary}

\passthrough{\lstinline!DiscoveryService.get\_capability\_summary()!}
returns a compact summary for task shaping and lightweight
orchestration.

Shape:

\begin{lstlisting}
{
  "roles": ["llm", "memory"],
  "capabilities_by_role": {
    "llm": ["chat", "summarize", "analyze"]
  },
  "modalities_by_role": {
    "llm": ["text"]
  },
  "deferred_available": false
}
\end{lstlisting}

\passthrough{\lstinline!deferred\_available!} is currently true when a
discovered backend has role \passthrough{\lstinline!cortex!}.

\subsection{13.8 Backend selection}\label{138-backend-selection}

Backend selection is handled by
\passthrough{\lstinline!BackendSelector!}.

A selection request can specify:

\begin{lstlisting}[language=Python]
@dataclass(slots=True)
class BackendSelectionRequest:
    role: str | None = None
    required_capabilities: list[str] | None = None
    required_modalities: list[str] | None = None
    runtime_preference: list[dict[str, Any]] | None = None
    healthy_only: bool = True
\end{lstlisting}

Selection filters backends in this order:

\begin{enumerate}
\def\labelenumi{\arabic{enumi}.}
\tightlist
\item
  \passthrough{\lstinline!healthy\_only!}
\item
  \passthrough{\lstinline!runtime\_preference!}
\item
  \passthrough{\lstinline!role!}
\item
  \passthrough{\lstinline!required\_capabilities!}
\item
  \passthrough{\lstinline!required\_modalities!}
\end{enumerate}

After filtering, candidates are sorted by:

\begin{lstlisting}
healthy first
priority descending
weight descending
name ascending
\end{lstlisting}

\subsection{13.9 Backend routing policy and
Planner}\label{139-backend-routing-policy-and-planner}

\passthrough{\lstinline!BackendRoutingPolicy!} performs primary healthy
selection first. If no healthy backend matches and degraded fallback is
allowed, it retries with \passthrough{\lstinline!healthy\_only=False!}
and may return a degraded backend.

\passthrough{\lstinline!Planner!} is intentionally thin. It converts
WorkPacket routing requirements into a backend selection request:

\begin{lstlisting}[language=Python]
class Planner:
    def select_backend(self, packet: WorkPacket):
        return self.routing_policy.select_backend(
            role=packet.required_role,
            required_capabilities=packet.required_capabilities,
            required_modalities=packet.required_modalities,
            runtime_preference=packet.routing_hints.runtime_preference,
        )
\end{lstlisting}

\begin{center}\rule{0.5\linewidth}{0.5pt}\end{center}

\section{RouterService and
BackendClient}\label{routerservice-and-backendclient}

\subsection{14.1 RouterService}\label{141-routerservice}

\passthrough{\lstinline!RouterService!} is the central dispatch layer
for WorkPackets.

It owns:

\begin{lstlisting}
backend selection
packet submission
accepted-work metadata
polling
batch send/retrieve operations
\end{lstlisting}

Dependencies:

\begin{lstlisting}
DiscoveryService
BackendClient
JobManager
Planner
\end{lstlisting}

\subsection{14.2 Sending one
WorkPacket}\label{142-sending-one-workpacket}

Main method:

\begin{lstlisting}[language=Python]
await router.send(packet)
\end{lstlisting}

Send flow:

\begin{enumerate}
\def\labelenumi{\arabic{enumi}.}
\tightlist
\item
  Resolve backend.
\item
  Submit packet through \passthrough{\lstinline!BackendClient!}.
\item
  If backend returns \passthrough{\lstinline!completed!} or
  \passthrough{\lstinline!failed!}, return that WorkResult.
\item
  If backend returns \passthrough{\lstinline!accepted!}, attach backend
  metadata and return accepted WorkResult.
\item
  If response shape is invalid, return failed WorkResult.
\end{enumerate}

Backend resolution first checks:

\begin{lstlisting}[language=Python]
packet.metadata["backend_name"]
\end{lstlisting}

If no explicit backend is found, RouterService calls:

\begin{lstlisting}[language=Python]
planner.select_backend(packet)
\end{lstlisting}

\subsection{14.3 Accepted WorkResult
contract}\label{143-accepted-workresult-contract}

When a backend accepts work asynchronously, RouterService returns a
\passthrough{\lstinline!WorkResult!} with:

\begin{lstlisting}
status = accepted
work_id = backend work id
backend_name = selected backend name
backend_model = selected backend model
metadata.backend_ref = serialized backend reference
metadata.work_details = work type and operation
\end{lstlisting}

\passthrough{\lstinline!backend\_ref!} contains only the information
needed to retrieve the result later:

\begin{lstlisting}
name
namespace
service_name
url
work_path
\end{lstlisting}

\passthrough{\lstinline!work\_details!} contains:

\begin{lstlisting}
work_type
operation
\end{lstlisting}

This is the contract for accepted or running work. Retrieval code should
use \passthrough{\lstinline!metadata["backend\_ref"]!}, not the older
\passthrough{\lstinline!backend\_desc!}.

\subsection{14.4 Retrieval and polling}\label{144-retrieval-and-polling}

Single retrieval:

\begin{lstlisting}[language=Python]
await router.retrieve_work(
    work_id=work_id,
    work_type=work_type,
    operation=operation,
    backend=backend_ref,
)
\end{lstlisting}

Blocking retrieval:

\begin{lstlisting}[language=Python]
await router.retrieve_completed_work(
    work_id=work_id,
    work_type=work_type,
    operation=operation,
    backend=backend_ref,
    poll_interval_seconds=1.0,
    timeout_seconds=120.0,
)
\end{lstlisting}

It polls while status is:

\begin{lstlisting}
accepted
running
\end{lstlisting}

and returns when the backend returns:

\begin{lstlisting}
completed
failed
\end{lstlisting}

If the timeout is reached, RouterService returns a failed WorkResult
with timeout metadata.

\subsection{14.5 Batch send/retrieve}\label{145-batch-sendretrieve}

Submit many WorkPackets:

\begin{lstlisting}[language=Python]
batch_id = await router.send_many(work_packets)
\end{lstlisting}

For each packet, RouterService starts a normal JobManager job:

\begin{lstlisting}
kind = router.sendmany
runner = execute_send
\end{lstlisting}

Retrieve many WorkResults:

\begin{lstlisting}[language=Python]
retrieve_batch_id = await router.retrieve_many(results)
\end{lstlisting}

For each result:

\begin{lstlisting}
completed or failed -> passthrough job
accepted or running -> retrieval job
\end{lstlisting}

Completed and failed results are not retrieved again. They are passed
through using:

\begin{lstlisting}
kind = router.retrievemany.passthrough
\end{lstlisting}

Accepted and running results use:

\begin{lstlisting}
kind = router.retrievemany
runner = execute_retrieve_completed
\end{lstlisting}

\passthrough{\lstinline!execute\_retrieve\_completed()!} requires:

\begin{lstlisting}
result.metadata["work_details"]
result.metadata["backend_ref"]
\end{lstlisting}

\subsection{14.6 BackendClient}\label{146-backendclient}

\passthrough{\lstinline!BackendClient!} is the internal HTTP client used
by RouterService and backend-scoped command handlers.

It supports both full \passthrough{\lstinline!BackendDescriptor!}
objects and serialized backend reference dictionaries.

Backend value access is normalized through:

\begin{lstlisting}[language=Python]
_backend_value(backend, key, default)
\end{lstlisting}

If \passthrough{\lstinline!backend!} is a dict, the value is read with
\passthrough{\lstinline!backend.get(key, default)!}. Otherwise it is
read with \passthrough{\lstinline!getattr(backend, key, default)!}.

Packet submission:

\begin{lstlisting}[language=Python]
await backend_client.submit_packet(
    backend=backend,
    packet=packet,
)
\end{lstlisting}

Work retrieval:

\begin{lstlisting}[language=Python]
await backend_client.retrieve_work(
    backend=backend_ref,
    work_id=work_id,
)
\end{lstlisting}

URL construction:

\begin{lstlisting}
<backend.url><backend.work_path>
<backend.url><backend.work_path>/<work_id>
\end{lstlisting}

The backend URL must start with \passthrough{\lstinline!http://!} or
\passthrough{\lstinline!https://!}.

Generic helpers:

\begin{lstlisting}[language=Python]
await backend_client.get_json(url=..., params=...)
await backend_client.post_json(url=..., payload=...)
\end{lstlisting}

These helpers are used for backend health checks, model status checks,
remote model loading, backend job status, and unload operations.

\begin{center}\rule{0.5\linewidth}{0.5pt}\end{center}

\section{Federation architecture}\label{federation-architecture}

Federation is the external membrane between Orin and the outside world.

External peers see only Federation-level concepts:

\begin{lstlisting}
FederationNetwork
NetworkMember
CapabilitySummary
FederatedWorkEnvelope
FederationWorkRecord
\end{lstlisting}

They must not see internal Orin implementation details.

\subsection{15.1 Source layout}\label{151-source-layout}

\begin{lstlisting}
src/federation/
  models.py
  settings.py
  storage.py
  network_registry.py
  join_tokens.py
  members.py
  identity.py
  envelopes.py
  policy.py
  app.py
  capabilities.py
  cortex_bridge.py
  federation_client.py
\end{lstlisting}

\subsection{15.2 Runtime composition}\label{152-runtime-composition}

During FastAPI lifespan startup, Federation:

\begin{enumerate}
\def\labelenumi{\arabic{enumi}.}
\tightlist
\item
  loads the federation configuration file
\item
  creates an internal HTTP client
\item
  creates an external HTTP client
\item
  creates \passthrough{\lstinline!MemoryFederationStorage!}
\item
  loads or creates local cluster identity
\item
  creates \passthrough{\lstinline!LocalNonceStore!}
\item
  creates \passthrough{\lstinline!NetworkRegistry!}
\item
  creates \passthrough{\lstinline!JoinTokenService!}
\item
  creates \passthrough{\lstinline!MemberService!}
\item
  creates \passthrough{\lstinline!CapabilityService!}
\item
  creates \passthrough{\lstinline!CortexBridge!}
\end{enumerate}

Persistent federation state is stored through Memory via
\passthrough{\lstinline!MemoryFederationStorage!}. Local files are used
for identity keys, nonce/cache state, and temporary local state.

\subsection{15.3 Settings}\label{153-settings}

\passthrough{\lstinline!FederationSettings!} is loaded from the shared
file-based configuration system:

\begin{lstlisting}[language=Python]
configuration.load_configuration_file("federation")
configuration.get_configuration("federation")
\end{lstlisting}

Main settings:

\begin{lstlisting}
app_name = "orin-federation"
protocol_version = "v1"
host = "0.0.0.0"
port = 8080
cluster_id: str | None = None
data_dir = /data/federation
private_key_path: Path | None = None
public_key_path: Path | None = None
memory_base_url = http://memory-service:8080
memory_work_path = /work
public_base_url: str | None = None
cortex_base_url = http://cortex-service:8080
cortex_work_path = /work
cortex_work_result_path = /work/{work_id}
cortex_network_path = /network
default_network_visibility = public
default_join_mode = token
request_ttl_seconds = 300
allowed_clock_skew_seconds = 60
nonce_ttl_seconds = 600
max_request_bytes = 1_000_000
enable_remote_work_submission = true
enable_capability_publish = true
enable_member_heartbeat = true
\end{lstlisting}

Configuration may provide key paths, but not raw key material. If key
paths are omitted, defaults are:

\begin{lstlisting}
/data/federation/identity/ed25519_private.key
/data/federation/identity/ed25519_public.key
\end{lstlisting}

\subsection{15.4 Persistent storage}\label{154-persistent-storage}

Federation does not own a separate SQLite database. Persistent
federation records are stored through the existing Memory service.

Supported record types:

\begin{lstlisting}
network
join_token
member
work
\end{lstlisting}

Expected Memory operations:

\begin{lstlisting}
federation_put_record
federation_get_record
federation_list_records
federation_delete_record
\end{lstlisting}

Storage keys:

\begin{lstlisting}
network:    network_id
join_token: token_id
member:     <network_id>:<cluster_id>
work:       <network_id>:<work_id>
\end{lstlisting}

\subsection{15.5 Local nonce store}\label{155-local-nonce-store}

\passthrough{\lstinline!LocalNonceStore!} is a local replay-protection
cache. It is intentionally local file state, not persistent Memory
state.

Default path:

\begin{lstlisting}
/data/federation/cache/nonces.json
\end{lstlisting}

Default TTL:

\begin{lstlisting}
600 seconds
\end{lstlisting}

\passthrough{\lstinline!check\_and\_remember(nonce)!} returns true when
the nonce is new and false when it has already been seen.

\subsection{15.6 Identity and
signatures}\label{156-identity-and-signatures}

Federation identity is Ed25519-based.

Public keys are encoded as:

\begin{lstlisting}
ed25519:<base64url>
\end{lstlisting}

Signatures are encoded as:

\begin{lstlisting}
ed25519:<base64url>
\end{lstlisting}

If \passthrough{\lstinline!settings.cluster\_id!} is not configured,
Federation derives a stable cluster ID from the public key:

\begin{lstlisting}
orin-<first_32_hex_chars_of_sha256(public_key_string)>
\end{lstlisting}

Canonical JSON for signatures:

\begin{lstlisting}[language=Python]
json.dumps(
    payload,
    sort_keys=True,
    separators=(",", ":"),
    ensure_ascii=False,
).encode("utf-8")
\end{lstlisting}

Private key files are written with mode \passthrough{\lstinline!0600!}.
Public key files are written with mode \passthrough{\lstinline!0644!}.

\subsection{15.7 Networks}\label{157-networks}

A \passthrough{\lstinline!FederationNetwork!} is the public concept that
Orin clusters can discover and join.

\begin{lstlisting}[language=Python]
class FederationNetwork(BaseModel):
    network_id: str
    slug: str
    name: str
    description: str | None = None
    visibility: "public | unlisted | private" = "public"
    join_mode: "token | approval | open | closed" = "token"
    owner_cluster_id: str
    policy: NetworkPolicy
    advertised_capabilities: list[CapabilitySummary]
    member_count: int
\end{lstlisting}

Slug rules:

\begin{lstlisting}
3-64 characters
lowercase letters
numbers
hyphens
must start and end with a letter or number
\end{lstlisting}

\passthrough{\lstinline!NetworkRegistry!} owns network creation, lookup,
listing, metadata updates, advertised capability updates, member-count
refresh, deletion, and public-safe views.

\subsection{15.8 Network policy}\label{158-network-policy}

\passthrough{\lstinline!NetworkPolicy!} is default-deny and allow-list
based.

Default allowed work types:

\begin{lstlisting}
llm
tool
\end{lstlisting}

Default allowed operations:

\begin{lstlisting}
chat
summarize
classify
extract
analyze
search
inspect
\end{lstlisting}

Policy fields include:

\begin{lstlisting}
allow_remote_work_submission
allow_remote_result_polling
allow_member_capability_publish
allowed_work_types
allowed_operations
max_payload_bytes
max_context_tokens
max_result_tokens
max_concurrent_jobs_per_member
expose_member_list
expose_exact_models
expose_runtime_metadata
\end{lstlisting}

There is no deny-list in the current model.

\subsection{15.9 Join tokens}\label{159-join-tokens}

Join tokens are bootstrap credentials only.

Raw token format:

\begin{lstlisting}
orin_join_<token_id>.<secret>
\end{lstlisting}

Stored token records contain only a hash:

\begin{lstlisting}
token_hash = sha256(raw_token)
\end{lstlisting}

The raw token is returned once and must never be persisted.

\passthrough{\lstinline!JoinTokenService!} owns token creation,
validation, redemption, revocation, expiry, deletion, and listing. It
does not create \passthrough{\lstinline!NetworkMember!} records.

\subsection{15.10 Members}\label{1510-members}

A \passthrough{\lstinline!NetworkMember!} is a joined cluster identity
inside a FederationNetwork.

\begin{lstlisting}[language=Python]
class NetworkMember(BaseModel):
    network_id: str
    cluster_id: str
    public_key: str
    role: "owner | admin | member | guest" = "member"
    status: "active | disabled | revoked | pending" = "active"
    allowed_work_types: list[str] = ["llm", "tool"]
    allowed_operations: list[str] = ["chat", "summarize", "analyze"]
    advertised_capabilities: list[CapabilitySummary]
\end{lstlisting}

Only active members can heartbeat, publish capabilities, or submit work.

A member can submit work only if:

\begin{lstlisting}
member.status == active
work_type in member.allowed_work_types
operation in member.allowed_operations
\end{lstlisting}

Both network policy and member policy must allow the requested work.

\subsection{15.11
FederatedWorkEnvelope}\label{1511-federatedworkenvelope}

External peers do not send raw internal WorkPackets directly. They send
a signed envelope:

\begin{lstlisting}[language=Python]
class FederatedWorkEnvelope(BaseModel):
    federation_version: str = "v1"
    network_id: str
    origin_cluster_id: str
    target_cluster_id: str | None = None
    request_id: str
    issued_at: datetime
    expires_at: datetime
    nonce: str
    packet: dict[str, Any]
    signature: str
    signature_algorithm: str = "ed25519"
    metadata: dict[str, Any] = Field(default_factory=dict)
\end{lstlisting}

The \passthrough{\lstinline!packet!} field is intentionally
\passthrough{\lstinline!dict[str, Any]!}, not a direct internal
\passthrough{\lstinline!WorkPacket!} type. This keeps the Federation
model layer decoupled from Cortex/common imports.

The envelope signature covers all envelope fields except
\passthrough{\lstinline!signature!}. The signed payload still includes
\passthrough{\lstinline!signature\_algorithm!}.

Envelope validation checks:

\begin{lstlisting}
federation version
signature algorithm
network_id, when expected
origin_cluster_id, when expected
issued_at / expires_at timing
signature
nonce replay, when nonce_store is provided
\end{lstlisting}

Envelope validation does not perform membership lookup, policy
enforcement, WorkPacket validation, quota checks, or routing.

\subsection{15.12 Policy checks and packet
sanitization}\label{1512-policy-checks-and-packet-sanitization}

Known runtime operations:

\begin{lstlisting}
chat
summarize
classify
extract
analyze
search
inspect
\end{lstlisting}

Federation policy checks include:

\begin{lstlisting}
remote work submission is enabled
task object exists
task.work_type exists and is a string
task.operation exists and is a string
operation is a known runtime operation
work type is allowed by network policy
operation is allowed by network policy
payload size is within network/app limits
requested context limit is within policy
requested result-token limit is within policy
member is active
member is allowed to submit the work type
member is allowed to submit the operation
\end{lstlisting}

Before a federated packet reaches Cortex, Federation removes dangerous
or trusted-looking metadata keys:

\begin{lstlisting}
backend_ref
backend_desc
selected_backend
internal
trusted
admin
command
deployment
kubernetes
ssh
\end{lstlisting}

Then it adds explicit federation-origin metadata:

\begin{lstlisting}[language=Python]
metadata["federation"] = {
    "network_id": network.network_id,
    "network_slug": network.slug,
    "origin_cluster_id": member.cluster_id,
    "member_role": member.role,
    "request_id": request_id,
}
\end{lstlisting}

Cortex should treat this as external-origin context, not trusted
internal state.

\subsection{15.13 Capability summaries}\label{1513-capability-summaries}

\passthrough{\lstinline!CapabilityService!} builds sanitized public
Federation capability summaries from local Cortex discovery state.

It reads:

\begin{lstlisting}
settings.cortex_network_url()
\end{lstlisting}

It accepts several Cortex network response shapes:

\begin{lstlisting}
[backend, backend]
{"result": [...]}
{"backends": [...]}
{"network": [...]}
{"items": [...]}
{"descriptors": [...]}
{"llm": [...], "tool": [...]}
\end{lstlisting}

Only public work types are advertised:

\begin{lstlisting}
llm
tool
\end{lstlisting}

Multiple backend descriptors are merged into one public summary per work
type so Federation does not reveal internal pod/service counts.

Capability summaries may include:

\begin{lstlisting}
work_type
operations
modalities
max_context_hint
max_result_tokens_hint
availability_hint
latency_hint
\end{lstlisting}

Exact model names are included only if
\passthrough{\lstinline!expose\_exact\_models=true!}.

Runtime metadata is included only if
\passthrough{\lstinline!expose\_runtime\_metadata=true!}, and only as
coarse safe hints.

\subsection{15.14 CortexBridge}\label{1514-cortexbridge}

\passthrough{\lstinline!CortexBridge!} bridges validated and sanitized
Federation work into local Cortex.

Expected Cortex contract:

\begin{lstlisting}
POST /work
GET  /work/{work_id}
\end{lstlisting}

Expected response shape:

\begin{lstlisting}
{
  "status": "accepted",
  "work_id": "...",
  "content": [],
  "backend_name": "...",
  "backend_model": "...",
  "error": null,
  "metadata": {}
}
\end{lstlisting}

Supported Cortex statuses:

\begin{lstlisting}
accepted
running
completed
failed
\end{lstlisting}

\passthrough{\lstinline!CortexBridge!} creates a
\passthrough{\lstinline!FederationWorkRecord!} and keeps separate IDs:

\begin{lstlisting}
origin_work_id
  Original WorkPacket id sent by the remote member.

cortex_work_id
  Internal Cortex job/work id used for local polling.

work_id
  Public Federation polling id returned to the remote member.
\end{lstlisting}

Current Phase 1 behavior: the public federation
\passthrough{\lstinline!work\_id!} is the Cortex work id. Later this can
become an independent \passthrough{\lstinline!fedwork:<id>!} value.

\subsection{15.15 FederationClient}\label{1515-federationclient}

\passthrough{\lstinline!FederationClient!} is the outbound client for
talking to another Federation node. It only talks to Federation
endpoints.

It must not talk directly to:

\begin{lstlisting}
Cortex
LLM
Memory
Tool
Kubernetes
backend services
\end{lstlisting}

Main methods:

\begin{lstlisting}
list_networks()
get_network()
get_network_capabilities()
join_network()
send_heartbeat()
publish_capabilities()
submit_work()
get_work_result()
submit_work_and_poll_once()
\end{lstlisting}

\subsection{15.16 Current Federation
flow}\label{1516-current-federation-flow}

Current supported flow:

\begin{enumerate}
\def\labelenumi{\arabic{enumi}.}
\tightlist
\item
  Create a FederationNetwork.
\item
  Create a join token.
\item
  Another cluster joins with token, cluster ID, and public key.
\item
  Member sends signed heartbeat.
\item
  Member publishes signed sanitized capabilities.
\item
  Member submits a signed FederatedWorkEnvelope.
\item
  Receiving Federation verifies envelope, membership, nonce, and policy.
\item
  Federation sanitizes the packet.
\item
  Federation submits the packet to local Cortex through CortexBridge.
\item
  Federation stores the FederationWorkRecord through Memory.
\item
  Remote member polls result through Federation.
\item
  Federation refreshes accepted/running records from Cortex and returns
  result or error.
\end{enumerate}

\subsection{15.17 Current Federation implementation
notes}\label{1517-current-federation-implementation-notes}

\begin{itemize}
\tightlist
\item
  Local/admin endpoints for creating networks and join tokens are
  currently unprotected. Later they should be restricted to the local
  command interface or explicit admin policy.
\item
  Quota and concurrency enforcement are represented in policy fields but
  are not fully implemented in the policy layer yet.
\item
  Federated backend descriptors for normal Cortex routing are still a
  later phase. Current remote work submission goes through explicit
  Federation endpoints and CortexBridge.
\item
  Federation settings intentionally do not include bootstrap peers in
  the current phase. Bootstrap/relay discovery belongs to a later phase.
\end{itemize}

\begin{center}\rule{0.5\linewidth}{0.5pt}\end{center}

\section{Primer and model execution}\label{primer-and-model-execution}

\passthrough{\lstinline!Primer!} is the shared runtime facade used by
Cortex and the LLM app.

It supports:

\begin{itemize}
\tightlist
\item
  backend loading: \passthrough{\lstinline!gguf!} or
  \passthrough{\lstinline!vllm!}
\item
  model provider loading: currently Hugging Face
\item
  model file download/ensure logic
\item
  model loading into backend
\item
  OpenAI-style message adapter loading
\item
  token counting
\item
  text chat
\item
  JSON chat
\item
  streaming text
\end{itemize}

Primer uses canonical \passthrough{\lstinline!RuntimeMessage!}
internally and renders messages through an adapter before sending them
to the backend.

\subsection{16.1 Message adapter}\label{161-message-adapter}

The current adapter is
\passthrough{\lstinline!OpenAIStyleMessageAdapter!}.

It renders:

\begin{itemize}
\tightlist
\item
  text parts into text
\item
  JSON parts into text
\item
  image parts into image blocks
\item
  unsupported parts into textual placeholders
\end{itemize}

When \passthrough{\lstinline!nothink=True!}, it injects:

\begin{lstlisting}
Do not output reasoning, chain-of-thought, or <think> tags. /no_think
\end{lstlisting}

If a model still emits \passthrough{\lstinline!<think>...</think>!}, the
adapter splits that into an internal reasoning message and a
user-visible final message.

\subsection{16.2 Generation policies}\label{162-generation-policies}

Current token policy:

\begin{lstlisting}
chat:      384
summarize: 640
classify:  64
extract:  128
analyze:  900
search:   256
inspect:  256
\end{lstlisting}

Current temperature policy:

\begin{lstlisting}
chat:      0.7
summarize: 0.3
classify:  0.1
extract:   0.1
analyze:   0.4
search:    0.2
inspect:   0.1
\end{lstlisting}

\passthrough{\lstinline!SAFETY\_TOKEN\_SIZE!} is currently
\passthrough{\lstinline!64!}.

\begin{center}\rule{0.5\linewidth}{0.5pt}\end{center}

\section{LLM backend lifecycle}\label{llm-backend-lifecycle}

Typical LLM lifecycle:

\begin{lstlisting}
POST /engine {"engine":"gguf"}
POST /load   {"model_id":"...", "repo_id":"...", "filename":"...", "tokenizer_id":"..."}
GET  /jobs/{job_id}
POST /work   WorkPacket
GET  /work/{work_id}
POST /unload
\end{lstlisting}

The LLM backend requires Primer readiness before accepting work. If
Primer is not ready, \passthrough{\lstinline!/work!} returns HTTP 409
with \passthrough{\lstinline!primer\_not\_ready!}.

LLM \passthrough{\lstinline!/model\_status!} returns useful routing
metadata after a model is loaded, including effective context size.
Discovery probing uses this to populate:

\begin{lstlisting}
BackendDescriptor.runtime.effective_n_ctx
\end{lstlisting}

\begin{center}\rule{0.5\linewidth}{0.5pt}\end{center}

\section{Memory backend}\label{memory-backend}

\subsection{18.1 Memory operations}\label{181-memory-operations}

Memory supports operations such as:

\begin{lstlisting}
create_summary
list_summaries
create_note
list_notes
create_memory_event
list_memory_events
create_memory_claim
list_memory_claims
search_memory_claims
retrieve
prompt_context
upsert_user
create_session
resolve_session
create_message
list_recent_messages
federation_put_record
federation_get_record
federation_list_records
federation_delete_record
\end{lstlisting}

The Federation operations are used by
\passthrough{\lstinline!MemoryFederationStorage!} for persistent
Federation state.

\subsection{18.2 Persistence}\label{182-persistence}

Default memory base directory:

\begin{lstlisting}
./cluster-memory
\end{lstlisting}

Default storage paths:

\begin{lstlisting}
cluster-memory/memory.db
cluster-memory/notes/
cluster-memory/summaries/
cluster-memory/docs/
cluster-memory/yaml/
\end{lstlisting}

\subsection{18.3 SQLite tables}\label{183-sqlite-tables}

Memory creates tables such as:

\begin{lstlisting}
users
chat_sessions
chat_messages
session_summaries
note_index
memory_events
memory_claims
work_items
\end{lstlisting}

\subsection{18.4 Prompt context}\label{184-prompt-context}

The \passthrough{\lstinline!prompt\_context!} operation returns:

\begin{lstlisting}
recent_messages
claims
summaries
retrieved_chunks
prompt_block
\end{lstlisting}

Cortex uses this for interpretation, shaping, last-assistant-message
retrieval, and final response building.

\subsection{18.5 Session behavior}\label{185-session-behavior}

\passthrough{\lstinline!resolve\_session!} uses inactivity-based
rollover:

\begin{enumerate}
\def\labelenumi{\arabic{enumi}.}
\tightlist
\item
  find latest active session for user/channel
\item
  if none exists, create a session
\item
  if the latest active session is older than
  \passthrough{\lstinline!max\_idle\_minutes!}, close it and create a
  new one
\item
  otherwise reuse the active session
\end{enumerate}

Default maximum idle time from Cortex is currently 60 minutes.

\begin{center}\rule{0.5\linewidth}{0.5pt}\end{center}

\section{Terminal client}\label{terminal-client}

\passthrough{\lstinline!terminal-chat.py!} is a lightweight local
terminal client for Cortex.

Default target:

\begin{lstlisting}
http://orin-gw:30080/work
\end{lstlisting}

Default user:

\begin{lstlisting}
terminal-user
\end{lstlisting}

The client sends canonical \passthrough{\lstinline!InferenceSession!}
payloads using \passthrough{\lstinline!RuntimeMessage!} content.

On startup it sends:

\begin{lstlisting}
/render
\end{lstlisting}

It separates command output from chat output by checking the response
header:

\begin{lstlisting}
Terminal-output: terminal
\end{lstlisting}

If that header is present, output goes into the command panel. Otherwise
output is treated as assistant chat.

For \passthrough{\lstinline!/attach!}, the terminal client should catch
Ctrl+C in streaming mode and return to the terminal loop instead of
exiting the whole client.

\begin{center}\rule{0.5\linewidth}{0.5pt}\end{center}

\section{Operations runbook}\label{operations-runbook}

\subsection{20.1 Bootstrap}\label{201-bootstrap}

Run the installer from the bootstrap/control host:

\begin{lstlisting}[language=bash]
./install-orin.sh
\end{lstlisting}

Useful environment overrides:

\begin{lstlisting}[language=bash]
ORIN_NAMESPACE=orin \
HOST_ALIAS=orin-gw \
HOST_NAME=orin-gw \
HOST_IP=<host-ip> \
PLATFORM=jetson \
NVIDIA_GPU=true \
REGISTRY=orin-gw:5000 \
IMAGE_VERSION=1.0.0 \
./install-orin.sh
\end{lstlisting}

\subsection{20.2 Health checks}\label{202-health-checks}

\begin{lstlisting}[language=bash]
kubectl get nodes -o wide
kubectl get pods -n orin -o wide
kubectl get svc -n orin
kubectl get deployments -n orin
\end{lstlisting}

Cortex:

\begin{lstlisting}[language=bash]
curl -s http://orin-gw:30080/health
curl -s http://orin-gw:30080/ready
curl -s http://orin-gw:30080/network
\end{lstlisting}

Memory inside cluster:

\begin{lstlisting}[language=bash]
kubectl exec -n orin netdebug -- sh -lc 'curl -s http://memory-service:8080/health'
\end{lstlisting}

\subsection{20.3 Debug pod}\label{203-debug-pod}

\begin{lstlisting}[language=bash]
kubectl run netdebug \
  -n orin \
  --image=nicolaka/netshoot \
  --restart=Never \
  --command -- /bin/sh -c "sleep 86400"
\end{lstlisting}

\subsection{20.4 Terminal client}\label{204-terminal-client}

\begin{lstlisting}[language=bash]
python terminal-chat.py
\end{lstlisting}

Initial command menu:

\begin{lstlisting}
/render
\end{lstlisting}

Useful command sequence:

\begin{lstlisting}
/commands
/cd Deployment
/list_nodes
/list_pods
\end{lstlisting}

\subsection{20.5 Add and install a
node}\label{205-add-and-install-a-node}

Add node to inventory:

\begin{lstlisting}
/add_node orin-node-1 --host orin-node-1 --ip 192.168.50.11 --platform jetson --gpu
\end{lstlisting}

Install k3s agent:

\begin{lstlisting}
/install_k3s orin-node-1
\end{lstlisting}

Sync labels and add role preference:

\begin{lstlisting}
/label_node orin-node-1 --role llm
\end{lstlisting}

List nodes:

\begin{lstlisting}
/list_nodes --verbose
\end{lstlisting}

\subsection{20.6 Deploy an LLM pod}\label{206-deploy-an-llm-pod}

\begin{lstlisting}
/deploy_pod llm light --platform jetson
\end{lstlisting}

or pin to a node:

\begin{lstlisting}
/deploy_pod llm light --node orin-node-1
\end{lstlisting}

List pods:

\begin{lstlisting}
/list_pods --verbose
\end{lstlisting}

\subsection{20.7 Load a model}\label{207-load-a-model}

Cortex-local model load from the Configuration folder:

\begin{lstlisting}
/cd /
/cd Configuration
/model_aliases
/load_model qwen35-4b
/job_status <job_id>
/attach
\end{lstlisting}

LLM backend model load from a selected backend folder:

\begin{lstlisting}
/cd /
/cd Cluster
/cd <llm-backend>
/model_status
/load_model qwen35-4b
/job_status <job_id>
/attach
\end{lstlisting}

For Hugging Face exploration:

\begin{lstlisting}
/cd /
/cd Configuration
/cd ModelDownloader
/huggingface search qwen
/huggingface gguf unsloth/Qwen3.5-9B-GGUF
\end{lstlisting}

\subsection{20.8 Federation manual test
outline}\label{208-federation-manual-test-outline}

A minimal Federation test flow:

\begin{lstlisting}
create network
create join token
join network with token, cluster_id, public_key
send signed heartbeat
publish signed capabilities
submit signed FederatedWorkEnvelope
poll federated work result
\end{lstlisting}

Exact curl examples should be documented separately once the
command-router Federation folder or test script is finalized.

\begin{center}\rule{0.5\linewidth}{0.5pt}\end{center}

\section{Known implementation notes and
cleanup}\label{known-implementation-notes-and-cleanup}

Current known cleanup items:

\begin{enumerate}
\def\labelenumi{\arabic{enumi}.}
\tightlist
\item
  The \passthrough{\lstinline!get\_job\_result!} branches in Harness
  should either be added to the lightweight interpretation model or
  removed.
\item
  \passthrough{\lstinline!RouterService.retrieve\_work()!} treats
  \passthrough{\lstinline!backend!} contractually as a dict but some
  error metadata still uses object-style
  \passthrough{\lstinline!getattr!}.
\item
  \passthrough{\lstinline!summarize\_jobs\_for\_prompt()!} includes
  payload details for completed jobs; this should become a compact
  result summary.
\item
  Some storage profile naming may still reference
  \passthrough{\lstinline!strong!} while the active pod template uses
  \passthrough{\lstinline!high!}.
\item
  Image naming should be cleaned up so \passthrough{\lstinline!amd64!}
  and \passthrough{\lstinline!arm64!} share the standard runtime image
  while Jetson keeps its own image.
\item
  \passthrough{\lstinline!PodImageConfig.image\_for()!} should treat
  \passthrough{\lstinline!arm64!} and \passthrough{\lstinline!amd64!} as
  separate platform values, not as one combined string.
\item
  \passthrough{\lstinline!DeploymentService.\_build\_pod\_factory\_defaults()!}
  should read registry/version from configuration instead of hardcoding
  defaults.
\item
  Routed streaming WorkPackets are not fully aligned if a backend
  returns a streaming response where RouterService expects JSON.
\item
  \passthrough{\lstinline!MemoryRuntime.handle\_egress()!} may remain a
  placeholder when memory work completes synchronously on POST.
\item
  Several memory error strings still say
  \passthrough{\lstinline!Summary Instance!} for other request types.
\item
  Mutable Pydantic defaults should consistently use
  \passthrough{\lstinline!Field(default\_factory=...)!}.
\item
  Direct \passthrough{\lstinline!/load!} for GGUF should consistently
  pass or infer the Hugging Face provider when
  \passthrough{\lstinline!repo\_id!} and
  \passthrough{\lstinline!filename!} are supplied.
\item
  Federation local/admin endpoints for network and token creation need
  local admin protection or command-router-only access.
\item
  Federation quota/concurrency fields exist in policy but are not fully
  enforced yet.
\item
  Federated backend descriptors for normal Cortex routing are a later
  phase.
\item
  Tool backend implementation remains future work.
\item
  Deployment inventory is still in-memory and should eventually be
  persisted.
\end{enumerate}

Removed stale items from older drafts:

\begin{lstlisting}
Federation as an empty/planned folder
old /task pipeline as the current background mechanism
memoryrequest field spelling
backend /network command as the current selected-backend display command
orin.ai/health_path, orin.ai/work_path, orin.ai/models_path as active discovery annotations
LLM /models as the active model-status discovery endpoint
accepted_not_executed Federation placeholder as current behavior
\end{lstlisting}

\begin{center}\rule{0.5\linewidth}{0.5pt}\end{center}

\section{Near-term architecture
direction}\label{near-term-architecture-direction}

Next architecture cleanup should focus on:

\begin{enumerate}
\def\labelenumi{\arabic{enumi}.}
\tightlist
\item
  keeping WorkPacket as the single execution unit across Cortex, LLM,
  Memory, Tool, and Federation paths
\item
  stabilizing Harness lightweight interpretation and background pipeline
  behavior
\item
  completing the standard-vs-Jetson image split
\item
  making service discovery metadata patching part of normal
  model/runtime operations
\item
  adding the first real Tool backend
\item
  protecting local/admin Federation endpoints
\item
  enforcing Federation quota/concurrency policy
\item
  adding command-router support for Federation operations
\item
  persisting deployment inventory beyond current in-memory inventory
\item
  replacing any remaining old gateway/orchestrator terminology with
  Cortex/LLM/Memory/Tool/Federation terminology
\end{enumerate}

The final direction is a public, no-pay overlay network where Orin
installs can discover FederationNetworks, join authorized networks with
tokens, and exchange signed policy-limited WorkPackets without exposing
internal cluster structure.

\end{document}
